% template-only:begin
% Template for a research chapter based on a paper. Copy it to chapters/<id>.tex (<id> as in
% research.json), replace <id> and <Research title>, delete this block, and add
% \include{chapters/<id>} to main.tex in research.json order. Write the text from the paper
% and projects/<id>/README.md; figures, tables and numbers come from projects/<id>/.
% template-only:end
\chapter{<Research title>}
\label{ch:<id>}

% --- Preface: the paper this chapter is based on (research.json `paper`) -------------------
% Uncomment and fill it in once the paper's entry is in reference.bib; delete it when the chapter
% has no paper. Keep it consistent with frontmatter/publications.tex. An item only the author can
% answer (contribution, permission) stays commented out, marked as a to-do, until it is answered.
% \begin{description}
%   \item[Publication] This chapter is based on \citet{<bibkey>}, <published in Journal (Year) |
%     accepted for publication in Journal | under review at Journal | in preparation>.
%   \item[Author's contribution] <What the author did (conception, data, analysis, writing) and
%     what the co-authors did.>
%   \item[Permission] <Reproduced with permission from the publisher | under the CC BY 4.0 licence.>
%   \item[Changes] <How this chapter differs from the published version, or "Formatting only.">
% \end{description}

\section{Introduction}
\label{sec:<id>-introduction}
% TODO: Narrow the general introduction (\cref{ch:introduction}) to this problem: the closest
% prior work, the gap and the objectives. Name the research questions this chapter answers
% (the RQ ids of \cref{sec:introduction-aims}).

\section{Methods}
\label{sec:<id>-methods}
% TODO: Data, methods and analysis in enough detail to repeat the study. Keep the paper's
% subsections (\subsection, labels subsec:<id>-...); describe a method shared with another
% chapter once and refer to it with \cref.

\section{Results}
\label{sec:<id>-results}
% TODO: What was found, without interpretation, with effect sizes and uncertainty. Figures,
% tables and numbers come from projects/<id>/ (examples at the end of this file).

\section{Discussion}
\label{sec:<id>-discussion}
% TODO: The principal findings, how they relate to prior work, and the limitations (also listed
% in projects/<id>/README.md).

\section{Summary}
\label{sec:<id>-summary}
% TODO: The answer to each research question of this chapter in a sentence, and how it leads to
% the next chapter. The general discussion (\cref{ch:conclusion}) builds on these summaries.

% --- Examples: figures, tables and numbers from projects/<id>/ -----------------------------
% Figure projects/<id>/figures/<name>.pdf, named without the extension:
% \begin{figure}[tbp]
%   \centering
%   \includegraphics[width=0.8\linewidth]{<id>/figures/<name>}
%   \caption{State the takeaway first, then the encodings and the uncertainty.}
%   \label{fig:<id>-<name>}
% \end{figure}
%
% Table: projects/<id>/tables/<name>.tex holds the tabular; caption and label stay here.
% \begin{table}[tbp]
%   \centering
%   \caption{What the table shows.}
%   \label{tab:<id>-<name>}
%   \input{../projects/<id>/tables/<name>}
% \end{table}
%
% Number defined in projects/<id>/values.tex (never type a result by hand):
%   We analysed \val{<id>-NSubjects} participants; the maximum depth was
%   \valqty{<id>-DepthMax}{\metre}.
